\section{Summary}
In conclusion, based on the analysis and testing conducted, the answers to our main questions are summarized as follows:

\subsection{1st Question}
Disney+ offers significantly less age-restricted content, supporting the claim that it is more suitable for children and families. While the platform predominantly caters to younger audiences, some exceptions exist, such as titles from the Marvel and Star Wars franchises, which appeal to older viewers as well. Overall, Disney+ maintains its position as the more child-friendly platform with a strong emphasis on family-oriented content.

\subsection{2nd Question}
On average, Disney+ provides content with slightly higher Rotten Tomatoes scores, indicating marginally better quality. However, the difference in average quality between Netflix and Disney+ is not substantial. Both platforms have a small portion of poorly rated titles, but the majority of their content holds above-average ratings.

While Disney+ edges out in quality, this is likely due to its smaller, more curated catalog. Netflix, with its larger and more diverse library, caters to a broader audience, offering a wide spectrum of content that includes both critically acclaimed and lower-rated titles. This variety highlights Netflix's strategy of balancing mainstream appeal with niche offerings, while Disney+ focuses on consistency. Importantly, both platforms demonstrate a commitment to delivering high-quality content, albeit through different strategies: Disney+ emphasizes quality through selectivity, whereas Netflix focuses on variety and inclusivity.
